\documentclass[paper=b5,fontsize=12pt]{jlreq}
\usepackage{amsmath, amssymb, mathcomp}
\usepackage{handwritenotebook}
\usepackage{mhchem}
\usepackage{chemfig}
\usepackage{pxrubrica}% ルビ
\usepackage{url,hyperref}
\usepackage{tikz, circuitikz, here}
\usetikzlibrary{snakes}
\usepackage{pict2e}

\begin{document}\randomspacing{
\section{目的}
精密に計量したものを用いて全量フラスコで純度100\unit\%の結晶を合成すること。
\section{理論}
まずは、布団から抽出した「やる気」0\unit{g}を精密に秤量する。続いて、ため息500\unit{mL}を全量フラスコに入れる。そこに「なんか熱っぽい気がする」という言い訳の飽和水溶液を標線まで注ぎ、静かに転倒混和する。この溶液を毛布でしっかりと包み込み、外部からの連絡を完全に遮断した状態で、約8時間ほど暗所にて静置する。すると、フラスコの底部に純度100\unit\%の「サボり」が美しく結晶化する。これが完璧な怠惰の合成法である。
}\end{document}
